% Modern editorial, markup and figure/layout contributions: CC-BY-SA-4.0.
% Historical text and illustrations: Leonhard Euler (1744), public domain.
% Component scope and upstream attribution: see RIGHTS.md.
\subsection*{Coroll. VI.}
\originalsection{22}\sourcepage{10}Porro ex iisdem substitutionibus erit, ut sequitur.
\[
\begin{aligned}
\text{Subtangens}&=\frac{y\,dx}{dy}=\frac{y}{p}\\
\text{Subnormalis}&=\frac{y\,dy}{dx}=py\\
\text{Tangens}&=\frac{y\,dw}{dy}=\frac{y\sqrt{(1+pp)}}{p}\\
\text{Normalis}&=\frac{y\,dw}{dx}=y\sqrt{(1+pp)}.
\end{aligned}
\]
Atque, pari modo, omnes quantitates finitæ ad curvam pertinentes, nisi
integralia involvant, per hujusmodi quantitates finitas ita exprimi poterunt,
ut nulla differentialia amplius inesse videantur.

\subsection*{Definitio IV.}
\originalsection{23}\textit{Maximi minimive Formula}, pro quovis Problemate, nobis erit ea quantitas, quæ in curva quæsita maximum minimumve valorem obtinere debet.

\subsection*{Coroll. I.}
\originalsection{24}Quoniam in omnibus Problematibus ad quæ hæc Methodus est accommodata, curva quæritur quæ, vel inter omnes, vel tantum inter innumeras curvas certo modo determinatas, maximi minimive proprietate gaudeat; hæc ipsa proprietas, quæ in curva quæsita maxima vel minima esse debet, erit quantitas; eaque exprimetur Formula, quam maximi minimive Formulam hic appellamus.

\subsection*{Coroll. II.}
\originalsection{25}\sourcepage{11}Cum autem maximi minimive proprietas ita proponi debeat, ut ad datam ac determinatam abscissam referatur; Formula maximi minimive quoque ad illam definitam abscissam debet referri.

\subsection*{Coroll. III.}
\originalsection{26}Erit igitur maximi minimive Formula, quantitas variabilis a longitudine abscissæ cujuscunque cui respondet pendens. Atque in quovis Problemate quæretur curva, pro qua, ad definitam abscissam, illa maximi minimive Formula maximum minimumve obtineat valorem.

\subsection*{Coroll. IV.}
\originalsection{27}Neque vero maximi minimive Formula a sola abscissa pendere potest: hoc enim si esset, pro omnibus curvis eidem abscissæ respondentibus eundem obtineret valorem, atque idcirco omnes æqualiter satisfacerent.

\subsection*{Coroll. V.}
\originalsection{28}Hanc ob rem quoque maximi minimive Formula, præter abscissam omnibus curvis quæ in considerationem veniunt communem, a qualibet curva peculiariter debet pendere; ita ut una sit, pro qua maximum minimumve valorem induere queat.

\subsection*{Scholion I.}
\originalsection{29}\marginpar{\footnotesize\hyperlink{fig:1}{Fig. 1.}}Quo hæc omnia clarius intelligantur, atque status Quæstionum in sequenti pertractandarum melius comprehendatur; ponamus, vel inter omnes omnino curvas, vel tantum inter innumerabiles certam quamdam proprietatem communem habentes, quæ eidem abscissæ $AZ$ respondeant, eam determinari debere,
\sourcepage{12}
pro qua valor formulæ $W$ sit maximus vel minimus. Ponamus huic Quæstioni satisfacere curvam $amz$, ita ut, quæcunque alia curva ad abscissam definitam $AZ$ referatur, valor formulæ $W$ vel fiat minor quam pro hac curva, vel major; prout in curva satisfaciente, $W$ vel maximum esse debet vel minimum. In hac igitur quæstione latissime patente, habemus primo abscissam determinatæ longitudinis $AZ$: deinde curva est quærenda, vel inter omnes omnino curvas ad eandem hanc abscissam relatas, vel tantum inter innumerabiles quibus una pluresve proprietates sint communes, prout quæstio ad methodum maximorum \& minimorum vel absolutam vel relativam est accommodata: tertio habemus eam quantitatem $W$, cujus valor in curva quæsita $amz$ maximus esse debet vel minimus; eritque igitur quantitas $W$ maximi minimive formula, sicut ea est definita. Nunc igitur statim apparet hanc formulam $W$ ita esse debere comparatam, ut ad omnes curvas quæ quidem concipi possunt accommodari queat. Primo scilicet a quantitate abscissæ definitæ $AZ$ debebit pendere; ita ut ea mutetur, valore ipsius $AZ$ mutato. Deinde etiam a natura cujusvis curvæ quæ quidem concipi potest peculiari modo debet pendere: nisi enim ita esset comparata, pro omnibus curvis eundem valorem sortiretur, quæstioque foret nulla. Quamobrem quantitas $W$, præter abscissam, in se quoque complecti debebit quantitates ad curvam ipsam pertinentes. Cum igitur omnis curva determinetur per relationem inter abscissam \& applicatam, quantitas $W$ debet esse conflata ex abscissa \& applicata, \& quantitatibus inde pendentibus. Hoc est, si abscissa indefinita ponatur $=x$, \& applicata respondens indefinita $=y$; quantitas $W$ esse debet functio binarum variabilium $x$ \& $y$. Quod cum ita sit, si curva quæcunque determinata concipiatur, atque ex ejus natura relatio inter $y$ \& $x$ in formula $W$ substituatur, ea definitum impetrabit valorem ad datam illam curvam atque ejus definitam abscissam pertinentem. Quoniam jam, pro aliis atque aliis curvis, formula $W$ diversos valores induit, etiamsi in omnibus abscissa eadem capiatur; manifestum est inter innumerabiles illas curvas \sourcepage{13}unam esse debere in qua valor formulæ $W$ maximus fiat vel minimus; atque ad hanc curvam pro data quacunque determinata quæstione inveniendam, Methodus tradenda est comparata.

\subsection*{Coroll. VI.}
\originalsection{30}Erit igitur maximi minimive formula $W$, functio quædam binarum variabilium $x$ \& $y$: quarum altera $x$ abscissam, altera $y$ applicatam denotat. In $W$ inesse igitur poterunt, non solum ipsæ variabiles $x$ \& $y$, sed etiam omnes quantitates ab iis pendentes, cujusmodi sunt $p$, $q$, $r$, $s$, \&c. quarum significationes supra tradidimus. Quinetiam formulæ integrales ex his ortæ quæcunque in $W$ inesse possunt: imo etiam debent; siquidem quæstio debeat esse determinata, uti mox ostendemus.

\subsection*{Coroll. VII.}
\originalsection{31}Proposita igitur ejusmodi formula $W$, seu functione ipsarum $x$ \& $y$, si quæstio ad methodum maximorum \& minimorum absolutam pertineat, ejusmodi æquatio inter $x$ \& $y$ desideratur, ut, si in $W$ valor ipsius $y$ per $x$ determinatus substituatur, atque ipsi $x$ valor definitus tribuatur; major prodeat quantitas pro $W$, vel minor, quam si ulla alia æquatio inter $x$ \& $y$ assumta fuisset.

\subsection*{Coroll. VIII.}
\originalsection{32}Hoc ergo pacto, quæstiones ad doctrinam linearum curvarum pertinentes ad Analysin puram revocari possunt. Atque vicissim, si hujus generis quæstio in Analysi pura sit proposita, ea ad doctrinam de lineis curvis poterit referri ac resolvi.

\subsection*{Scholion II.}
\originalsection{33}Quanquam hujus generis quæstiones ad puram Analysin \sourcepage{14}reduci possunt, tamen expedit eas cum doctrina linearum curvarum conjungere. Quod si enim animum a lineis curvis abducere, atque ad solas quantitates absolutas firmare velimus; quæstiones primum ipsæ admodum fierent abstrusæ \& inelegantes, ususque earum ac dignitas minus conspiceretur: Deinde etiam methodus resolvendi hujusmodi quæstiones, si in solis quantitatibus abstractis proponeretur, nimium foret abstrusa \& molesta; cum tamen eadem, per inspectionem figurarum \& quantitatum repræsentationem linearem, mirifice adjuvetur atque intellectu facilis reddatur. Hanc ob causam, etsi hujus generis quæstiones, cum ad quantitates abstractas, tum concretas applicari possunt, tamen eas ad lineas curvas commodissime traducemus \& resolvemus. Scilicet quoties æquatio ejusmodi inter $x$ \& $y$ quæritur, ut formula quædam proposita \& composita ex $x$ \& $y$, si ex illa æquatione quæsita valor ipsius $y$ subrogetur, \& ipsi $x$ determinatus valor tribuatur, maxima fiat vel minima: tum semper quæstionem transferemus ad inventionem lineæ curvæ, cujus abscissa sit $x$, \& applicata $y$, pro qua illa formula $W$ fiat maxima vel minima, si abscissa $x$ datæ magnitudinis capiatur. His igitur notatis, natura hujusmodi quæstionum satis luculenter perspicitur: nisi forte cuiquam adhuc dubium creat ambigua locutio de maximo \& minimo simul. Verum ne hic quidem ulla adest ambiguitas; nam etsi methodus ipsa æque monstrat maxima \& minima, tamen in quovis casu facile erit discernere, utrum solutio præbeat maximum an minimum. Sæpe numero autem evenire potest, ut in data quæstione tam maximum quam minimum locum obtineat, atque his casibus solutio erit duplex, altera monstrante maximum, altera minimum. Plerumque autem alterutrum, scilicet vel maximum vel minimum solet esse impossibile; quod evenit, si maximi minimive formula in infinitum vel crescere, vel decrescere potest; his enim casibus, vel non dabitur maximum, vel non minimum. Usu venire etiam potest, ut formula proposita $W$ in infinitum tam crescere quam decrescere queat, atque his casibus nulla prorsus solutio locum \sourcepage{15}habebit. Hæc autem discrimina cuncta ipse calculus post solutionem perpetuo monstrabit.

\subsection*{Propositio I. Theorema.}
\originalsection{34}\textit{Ut per maximi minimive formulam $W$ curva determinetur $amz$, quæ præ omnibus reliquis satisfaciat, formula $W$ debet esse quantitas integralis indefinita, quæ, nisi data assumatur relatio inter $x$ \& $y$, integrari nequeat.}

\subsection*{Demonstratio.}
Ponamus enim formulam $W$ integralia indefinita non involvere; erit ea functio quantitatum $x$ \& $y$, indeque pendentium $p$, $q$, $r$, $s$, \&c. vel algebraica, vel talis transcendens quæ sine assumta relatione inter $x$ \& $y$ exhiberi possit; quod evenit, si vel logarithmi harum quantitatum, vel arcus circulares, vel aliæ hujusmodi quantitates transcendentes definitæ ingrediantur, quæ algebraicis æquivalentes sunt censendæ. Quod si jam $W$ ponatur functio talis ipsarum $x$ \& $y$ tantum, manifestum est valorem formulæ $W$, quem pro data curva $amz$ ad datam abscissam $AZ$ relata obtinet, tantum ab ultima applicata $Zz$ pendere; atque pro omnibus curvis in $Z$ eandem applicatam $Zz$ habentibus fore eundem; atque adeo tali formula $W$ indoles totius curvæ non determinabitur, sed tantum positio extremi ejus puncti $z$; si in $W$ præter $x$ \& $y$ etiam quantitas $p$ insit, tum præter longitudinem applicatæ $Zz$ positio tangentis curvæ in $z$, seu positio ultimi elementi in $z$ determinabitur. Sin autem insuper $q$ ingrediatur, tum positio binorum elementorum curvæ contiguorum in $z$ determinabitur, \& ita porro. Ex quibus sequitur, si fuerit $W$ functio determinata ipsarum $x$, $y$, $p$, $q$, $r$, \&c. tum per illam tantum curvæ portionem infinite parvam circa extremitatem $z$ determinari; atque pro omnibus curvis in eandem extremitatem desinentibus eundem valorem ipsius $W$ esse proditurum. Ut itaque per formulam $W$ tota curva $amz$, quatenus toti abscissæ $AZ$ respondet, definiatur, \sourcepage{16}formulam $W$ ita oportet esse comparatam, ut ejus valor ad determinatam curvam $amz$ applicatus, a positione singulorum elementorum hujus curvæ intra terminos $a$ \& $z$ sitorum pendeat. Hoc autem evenire non potest, nisi quantitas $W$ sit formula integralis indefinita, quæ generatim sine assumta æquatione inter $x$ \& $y$ integrationem non admittat. Q. E. D.

\subsection*{Coroll. I.}
\originalsection{35}Nisi igitur maximi minimive formula $W$ sit quantitas integralis indefinita, nequidem linea curva in qua valor ipsius $W$ sit maximus vel minimus determinabitur; atque adeo quæstio de invenienda curva, in qua esset $W$ maximum vel minimum erit nulla.

\subsection*{Coroll. II.}
\originalsection{36}Ut igitur curva assignari possit, in qua valor ipsius $W$ præ aliis sit maximus vel minimus, formula $W$ talem formam $\int Z\,dx$ habere debet; atque quantitatem $Z$ ita comparatam esse oportet ut differentiale $Z\,dx$, nisi æquatio statuatur inter $x$ \& $y$, integrari nequeat.

\subsection*{Scholion.}
\originalsection{37}Quoniam maximi minimive formula $W$ debet esse integrale formulæ differentialis indefinitæ primi gradus; hoc est cujus integrale fiat quantitas finita; ea formula differentialis semper ad hujusmodi formam $Z\,dx$ poterit reduci, ope litterarum $p$, $q$, $r$, \&c. Et hanc ob rem in sequentibus maximi minimive formula perpetuo per $\int Z\,dx$ nobis indicabitur. Erit autem $Z$ functio non solum quantitatum $x$ \& $y$, sed etiam continebit litteras $p$, $q$, $r$, \&c. Ita si area $AazZ$ debeat esse maxima vel minima, formula $W$ abibit in $\int y\,dx$; \& si superficies solidi rotundi quod generatur rotatione curvæ $amz$ circa axem $AZ$ debeat esse maxima vel minima, erit $W=\int y\,dx\sqrt{(1+pp)}$;
